%------------------------------------------------------------------------------%
\section{Teste de Divergência}

A partir desse ponto começamos a construir um conjunto de ferramentas
que nos ajudam a decidir se uma série converge ou não.
Note que, por enquanto, não estamos preocupados com a soma da série,
queremos apenas saber se a soma existe.
O primeiro desses resultados, que chamamos de
\emph{Teste de Divergência} ou
\emph{Teste do Termo Geral},
nos ajuda a identificar muitas séries que não podem ser
convergentes com relativa facilidade .
Dessa forma, sempre o aplicamos primeiro para descartar rapidamente essas séries.
O teorema a seguir é o principal resultado que nos permite
construir o teste da divergência.

%------------------------------------------------------------------------------%
\begin{teorema}{Limite do Termo Geral}{teo:limite-termo-geral}
  Se a série
  \(
    \sum_{n=1}^\infty a_n
  \)
  é convergente $a_n$ tende para zero.
\end{teorema}
%------------------------------------------------------------------------------%
\begin{proof}
  Como a série é convergente, a sequência de somas parciais
  $S_n$ é convergente, ou seja,
  \[
    \lim S_n=\lim S_{n-1}=L
  \]
  observando que $S_n-S_{n-1}=a_n$ concluirmos que
  \[
    \lim a_n
    = \lim (S_n-S_{n-1})
    = L - L
    = 0              \,.
  \]
\end{proof}
%------------------------------------------------------------------------------%

O teste de divergência é uma consequência direta desse teorema.

%------------------------------------------------------------------------------%
\begin{teorema}{Teste de Divergência}{teo:teste-divergencia}
  Se\index{Teste!divergência}\index{Teste!termo geral}
  a sequência $a_n$ não converge para zero, então a série
  \(
    \sum_{n=1}^\infty a_n
  \)
  diverge.
\end{teorema}
%------------------------------------------------------------------------------%

Uma forma para provar que uma série diverge é mostrar que o seu termo geral $a_n$
ou converge para um limite diferente de zero ou diverge. Mas tome cuidado, isso não quer
dizer que toda série que diverge satisfaz a propriedade que o termo geral não converge
para zero, ou, analogamente, não é verdade que toda série cujo termo geral converge para
zero é uma série convergente. Os próximos exemplos ilustram o uso desse teste em
algumas séries específicas.

%------------------------------------------------------------------------------%
\begin{example}
  Mostrar que a série
  \(\;
    \sum_{n=1}^\infty ar^{n-1}
  \;\)
  diverge quando $\abs{r}\geq 1$.
  \solution
  A série geométrica
  \[
    \sum_{n=1}^\infty ar^{n-1}
  \]
  diverge quando $\abs{r}\geq 1$, pois, nesse caso, $ar^{n-1}$ não converge para zero.
  Verificamos isso observando que se $r=1$
  \[
    ar^n \to a \neq 0
  \]
  enquanto que, se $\abs{r}>1$ ou $r=-1$, $ar^{n-1}$ diverge.
\end{example}
%------------------------------------------------------------------------------%

%------------------------------------------------------------------------------%
\begin{example}
  Verificar a convergência da série
  \(\;
    \sum_{n=0}^\infty \left(\sqrt{2}\right)^n
  \)
  \solution
  A série
  \(
    \sum_{n=0}^\infty \left(\sqrt{2}\right)^n
  \)
  diverge, uma vez que
  \(
    \lim \left(\sqrt{2}\right)^n\!\!=\infty
  \)
  pois
  \(
    \sqrt{2}>1
  \).
\end{example}
%------------------------------------------------------------------------------%

%------------------------------------------------------------------------------%
\begin{example}
  Verificar a convergência da série
  \(\;
    \sum_{n=0}^\infty \cos(n\pi)
  \)
  \solution
  Essa série diverge, pois
  \(
    \lim \cos(n\pi)
  \)
  não existe, já que
  \[
    \cos(n\pi)=(-1)^n
  \]
  que alterna entre $-1$ e $1$.
\end{example}
%------------------------------------------------------------------------------%

%------------------------------------------------------------------------------%
\begin{example}
  Verificar a convergência da série
  \(\;
    \sum_{n=1}^\infty \frac{n^n}{n!}   \,.
  \)
  \solution
  A série
  \(
    \sum_{n=1}^\infty \frac{n^n}{n!}
  \)
  diverge, pois
  \(
    \lim \frac{n^n}{n!}=\infty
  \).
  De fato, note que, se $n\geq 2$,
  \[ n! \;=\;     1\cdot 2\cdot 3\cdots n
        \;\leq\;  \underbrace{n\cdot n \cdots n}_{n-1 \text{ fatores }}
        \;=\; n^{(n-1)}
  \]
  Então, multiplicando ambos os lados por $n$
  \[
    n\,n!\leq n^{n} \qquad\Rightarrow\qquad n\leq \frac{n^n}{n!}
  \]
  Como $n\to\infty$, seque que
  \(
    \frac{\,n^n\,}{n!}\to \infty
  \).
\end{example}
%------------------------------------------------------------------------------%

%------------------------------------------------------------------------------%
\begin{example}
  Verificar a convergência da série
  \(\;
    \sum_{n=1}^\infty \ln\left(\frac{n}{2n+1}\right)
  \).
  \solution
  A série \(\sum_{n=1}^\infty \ln\left(\frac{n}{2n+1}\right)\) diverge, pois
  \[
    \lim\ln\left(\frac{n}{2n+1}\right) \;=\; \ln\left(\lim \frac{n}{2n+1}\right)
                                       \;=\; \ln{\frac{1}{2}}
                                       \;\neq\; 0     \,.
  \]
\end{example}
%------------------------------------------------------------------------------%

A seguir apresentamos um exemplo clássico de uma série cujo termo geral converge para
zero mas a série diverge. Primeiro vamos apresentar a definição dessa
série e em seguida verificamos que ela é divergente.

%------------------------------------------------------------------------------%
\begin{definicao}{Série Harmônica}{def:serie-harmonica}
  A \emph{Série Harmônica}\index{Série!harmônica} é a série
  \(\quad
    \sum_{n=1}^\infty \frac{1}{\;n\;}
  \).
\end{definicao}
%------------------------------------------------------------------------------%

%------------------------------------------------------------------------------%
\begin{example}
  Verificar a convergência da série harmônica.
  \solution
  Para ver que essa série diverge, observe que
  \begin{align*}
   	\!\!\!S_{2^n} & = 1 + \frac{1}{2} + \cdots + \frac{1}{2^n} \\[3mm]
  	              & = 1 + \frac{1}{2}
                    + \left( \frac{1}{3} + \frac{1}{4} \right)
                    + \left( \frac{1}{5} + \frac{1}{6} + \frac{1}{7} + \frac{1}{8} \right)
                    + \cdots
                    + \left( \frac{1}{2^{n-1}+1} + \cdots + \frac{1}{2^n} \right) \\[3mm]
  	              & > 1 + \frac{1}{2}
                    + \left( \frac{1}{3} + \frac{1}{4} \right)
                    + \left( \frac{1}{8} + \frac{1}{8} + \frac{1}{8} + \frac{1}{8} \right)
                    + \cdots
                    + \left( \frac{1}{2^n} + \cdots + \frac{1}{2^n} \right) \\[3mm]
  	              & = 1 + \frac{1}{2}
                    + \frac{2}{4}
                    + \frac{4}{8} + \cdots
                    + \frac{2^{n-1}}{2^n}  \\[3mm]
  	              & = 1 + \frac{n}{2} \to \infty    \,.
  \end{align*}
  Portanto a série diverge, mesmo que $\sfrac{1}{n}\to0$.
\end{example}
%------------------------------------------------------------------------------%

%------------------------------------------------------------------------------%
